\documentclass[11pt]{article}
\usepackage{preamble}
\setlength{\parskip}{4pt}

\begin{document}

\daybanner{AP Statistics \textperiodcentered\ Topic 5.3 (CED 2.12) \textperiodcentered\ B: 2026-11-13 \textperiodcentered\ E: 2026-12-09 \textperiodcentered\ TEACHER KEY}{When Is the Mean Nearly Normal?}

\sectionbanner{CED TETHER}

{\small
\begin{itemize}[leftmargin=1.5em,itemsep=2pt,topsep=2pt]
  \item \textbf{Skill 3.C} --- Describe probability distributions.
  \item \textbf{Enduring Understanding UNC-3} --- Probabilistic reasoning allows us to anticipate patterns in data.
  \item \textbf{LO UNC-3.H} --- Estimate sampling distributions using simulation. [Skill 3.C]
  \item \textbf{EK UNC-3.H.1} --- A sampling distribution of a statistic is the distribution of values for the statistic for all possible samples of a given size from a given population.
  \item \textbf{EK UNC-3.H.2} --- The central limit theorem (CLT) states that when the sample size is sufficiently large, a sampling distribution of the mean of a random variable will be approximately normally distributed.
  \item \textbf{EK UNC-3.H.3} --- The central limit theorem requires that the sample values are independent of each other and that n is sufficiently large.
  \item \textbf{EK UNC-3.H.4} --- A randomization distribution is a collection of statistics generated by simulation assuming known values for the parameters.
  \item \textbf{EK UNC-3.H.5} --- The sampling distribution of a statistic can be simulated by generating repeated random samples from a population.
\end{itemize}
}
\textbf{Essential question:} How do population shape and sample size determine whether a normal model is trustworthy for sample means?\par
\textbf{DOK-3 skill:} 4.B \textperiodcentered\ \textbf{FRQ pattern:} {\small\texttt{clt-\allowbreak{}when-\allowbreak{}does-\allowbreak{}normality-\allowbreak{}apply}}\par
\textbf{DOK rationale:} Part (c) coordinates population shape, sample size, independence, and sampling-distribution parameters to adjudicate two normality claims and bound which probability model a reader may use.\par

\frameworkphaseheader{Do Now (first take) $\rightarrow$ Explore (video + follow-along) $\rightarrow$ Exit (finish + turn in)}{1 $\rightarrow$ 3}{45}{%
\begin{itemize}[leftmargin=1.5em,itemsep=2pt,topsep=2pt]
  \item Board slide up at the bell. Ask for a sample-size choice before displaying a rule of thumb.
  \item Begin the 5.3 video follow-along at minute 5.
  \item Return to the board slide at minute 33. Circulate and require students to name which distribution becomes approximately normal.
  \item Collect at the door. Score part (c) only, E/P/I.
\end{itemize}
}{%
\begin{itemize}[leftmargin=1.5em,itemsep=2pt,topsep=2pt]
  \item Choose n equals 4 or n equals 40 and cite population shape or sample size.
  \item Complete the follow-along about simulated sampling distributions, the CLT, and randomization distributions.
  \item Finish (a)--(c), revise the first take in the exit line, and turn in the sheet.
\end{itemize}
}{%
\begin{itemize}[leftmargin=1.5em,itemsep=2pt,topsep=2pt]
  \item Does averaging four values erase a long right tail automatically?
  \item Which distribution does the CLT describe: commuters or sample means?
  \item What happens to $15/\sqrt{n}$ when $n$ grows from 4 to 40?
  \item Which probability calculation becomes questionable if the n equals 4 shape is not normal?
\end{itemize}
}{Solo teacher. Circulate ~5 pairs. Require shape, sample size, and independence in every normality defense; redirect claims that the population itself changes shape.}

\begin{callout}{\IconWarn\ WATCH FOR}{calloutred}
\begin{itemize}[leftmargin=1.5em,itemsep=2pt,topsep=2pt]
  \item Invoking the CLT for n equals 4 without considering the strong population skew.
  \item Saying forty individual commute times have a normal distribution.
  \item Using 15 minutes, rather than $15/\sqrt{n}$, for the spread of sample means.
\end{itemize}
\end{callout}

\sectionbanner{SCORING PART (c) --- E / P / I}

\textbf{Expected elements}
\begin{itemize}[leftmargin=1.5em,itemsep=2pt,topsep=2pt]
  \item \textbf{rejects-small-n-normality} (required) --- Rejects the automatic normal claim for n equals 4 because the source population is strongly right-skewed
  \item \textbf{supports-large-n-normality} (required) --- Uses independence and sufficiently large n equals 40 to invoke the central limit theorem
  \item \textbf{uses-large-n-spread} (required) --- States that the n equals 40 sampling distribution is centered at 25 minutes with standard deviation about 2.37 minutes
  \item \textbf{distinguishes-distributions} (required) --- Clarifies that the CLT concerns the distribution of sample means, not the distribution of individual commute times
  \item \textbf{states-model-consequence} (required) --- Explains that a normal approximation at n equals 4 could give inaccurate tail probabilities
\end{itemize}

{\small\renewcommand{\arraystretch}{1.25}
\noindent\begin{tabular}{|p{0.5in}|p{5.9in}|}
\hline
\textbf{E} & Rejects normality at n equals 4, supports it at n equals 40 with independence and the CLT, uses the 2.37-minute standard deviation, distinguishes the two distributions, and explains the tail-probability consequence. \\ \hline
\textbf{P} & Chooses the appropriate sample sizes and mentions skew and the CLT but omits independence, the computed spread, the distribution distinction, or the consequence for probability calculations. \\ \hline
\textbf{I} & Says every average is normal, says the population itself becomes normal at n equals 40, or uses 15 minutes as the sampling-distribution standard deviation for both sample sizes. \\ \hline
\end{tabular}}

\textbf{Common mistakes}
\begin{itemize}[leftmargin=1.5em,itemsep=2pt,topsep=2pt]
  \item Treating the central limit theorem as applying at every sample size
  \item Saying the individual commute-time population becomes normal when n increases
  \item Dividing the mean, rather than the standard deviation, by $\sqrt{n}$
  \item Using the number of repeated samples instead of the observations per sample as $n$
\end{itemize}

\clearpage
\focusbanner{TODAY'S PROBLEM --- annotated}

Individual commute times in a very large city are strongly right-skewed, with population mean $\mu=25$ minutes and standard deviation $\sigma=15$ minutes. Consider independent random samples.\par\smallskip \textbf{Rae:} ``Any average is approximately normal, even for $n=4$.''\par \textbf{Sol:} ``With this long right tail, $n=4$ may be too small; $n=40$ is more defensible.''\par
\begin{center}
\begin{tikzpicture}
\begin{axis}[
  width=0.44\linewidth,
  height=1.20in,
  ybar interval, ymin=0, ymax=32,
  xmin=0, xmax=80, xtick={0,10,20,30,40,50,60,70,80},
  xlabel={Individual commute time (minutes)}, ylabel={Relative frequency (percent)},
  ymajorgrids, tick label style={font=\small}, label style={font=\small},
  bar shift=0pt, x tick label as interval=false
]
\addplot[fill=calloutblue, draw=dayblue] coordinates {
  (0,12)
  (10,30)
  (20,28)
  (30,15)
  (40,8)
  (50,4)
  (60,2)
  (70,1)
  (80,0)
};
\end{axis}
\end{tikzpicture}
\end{center}
\begin{firsttakebox}{FIRST TAKE (5 min)}
For which sample size, $n=4$ or $n=40$, would you trust a normal model for the sample mean? Name one feature behind your choice.\par
\answer{Not graded. Expect `averages are normal' as an initial rule; the revision should add population shape and sufficiently large n.}
\end{firsttakebox}

\textit{Rules callout ``CLT: SHAPE, SIZE, AND THE RIGHT DISTRIBUTION (keep on the page)'' is printed on the student sheet.}\par\medskip

\dokbadge{1}~\textbf{(a)} Identify the population shape and state the population mean and standard deviation with units.\par
\answer{The population distribution of individual commute times is strongly right-skewed. Its mean is $\mu=25$ minutes and its standard deviation is $\sigma=15$ minutes.}

\dokbadge{2}~\textbf{(b)} For $n=4$ and $n=40$, calculate the mean and standard deviation of the sampling distribution of $\bar{x}$.\par
\answer{For either sample size, $\mu_{\bar{x}}=\mu=25$ minutes. For $n=4$, $\sigma_{\bar{x}}=15/\sqrt{4}=7.5$ minutes. For $n=40$, $\sigma_{\bar{x}}=15/\sqrt{40}\approx2.3717$ minutes, or about $2.37$ minutes.}

\dokbadge{3}~\textbf{(c)} Adjudicate both claims for $n=4$ and $n=40$. Cite the histogram shape, independence, and the CLT's sample-size requirement; use the $n=40$ standard deviation; and explain one consequence of ignoring skew when using a normal model at $n=4$.\par
\answer{Rae's claim is not warranted at $n=4$: averaging only four observations from a strongly right-skewed population need not remove the skew. Sol is warranted at $n=40$: independence and the large sample size let the CLT support an approximately normal sampling distribution centered at 25 minutes with standard deviation $15/\sqrt{40}\approx2.37$ minutes. The CLT concerns sample means, not individual commute times. Ignoring the population skew at $n=4$ could make normal-curve tail probabilities inaccurate, making rare sample means look too common or too rare.}

\begin{summaryexitbox}{EXIT}
One thing the video changed about my first take: \hrulefill
\end{summaryexitbox}

\end{document}
